A detector that scores 99\% on a machine-generated-text benchmark may have modelled how machine language differs from human language, or how one corpus happened to be punctuated and encoded, and the headline number does not distinguish the two. We propose a decomposition that measures the difference. A surface-only model reads 47 orthographic features and never sees a word, a content-only model reads text with punctuation, casing and non-ASCII stripped, and a full model is a fine-tuned transformer on raw text. Across DAIGT V2 and HC3, five model families and eight configurations per corpus, the two benchmarks separate. On HC3 the surface and content arms are indistinguishable under a paired test, 3.20\% against 3.26\% error, and the null holds on all five of five independent partitions with the sign changing between them. On DAIGT V2 content is stronger by a factor of 7.9 in error. Closing the document-length channel both arms share widens the DAIGT V2 gap to 10.6 times and turns HC3's parity into a 2.91-point advantage for orthography. A SHAP attribution locates the asymmetry, since one feature, the rate of spaces preceding punctuation, dominates HC3 at 2.3 times the next while no orthographic feature dominates DAIGT V2. Three controls change what the remaining results mean. BERT emits identical token identifiers for strings differing only by that whitespace cue, so it cannot represent the cue and still reaches 0.9916 weighted F1, making the cue sufficient in isolation but unnecessary in practice. A pipeline control confirms this null is not a cleaning step that failed to execute. A label-free control shows these models assign the human label to uniform random character strings in 100\% of cases at 0.98 mean confidence, so an adversarial collapse to 0.3644 weighted F1 cannot be read as a targeted vulnerability. Surface form is informative on both corpora, and we claim only that it reaches parity with content on one.
